%%%%%%%%%%%%%%%%%%%%%%%%%%%%%%%%%%%%%%%%%%%%%%%%%%%%%%%%%%%%%%%%%%%%%%%%%%%%%%%%
%                                                                              %
% BlackArch Linux Guide                                                        %
%                                                                              %
%%%%%%%%%%%%%%%%%%%%%%%%%%%%%%%%%%%%%%%%%%%%%%%%%%%%%%%%%%%%%%%%%%%%%%%%%%%%%%%%

\documentclass[a4paper, oneside, 11pt]{book}

%%% INCLUDES %%%
\renewcommand{\familydefault}{\sfdefault}

\usepackage{array}
\usepackage{color}
\usepackage{enumerate}
\usepackage{fancyhdr}
\usepackage{fancyvrb}
\usepackage{geometry}
\usepackage{graphicx}
\usepackage{html}
\usepackage{hyperref}
\usepackage{ifpdf}
\usepackage{listings}
\usepackage{pstricks}
\usepackage{supertabular}
\usepackage{tocloft}
\usepackage[utf8]{inputenc}

%%% LAYOUT %%%
\setlength{\parindent}{0em}
\setlength{\parskip}{1.5ex plus0.5ex minus0.5ex}
\geometry{left=2.5cm, textwidth=16cm, top=3cm, textheight=25cm, bottom=3cm}
\widowpenalty=2000
\clubpenalty=1000
\frenchspacing
\sloppy
\pagecolor[HTML]{FFFFFF}
\color[HTML]{333333}
\setcounter{tocdepth}{10}
\setcounter{secnumdepth}{10}

\hypersetup{
  pdftitle={BlackArch Linux, The BlackArch Linux Guide},
  pdfsubject={BlackArch Linux, The BlackArch Linux Guide},
  pdfauthor={BlackArch Linux, BlackArch Linux},
  pdfkeywords={BlackArch Linux, Penetration Testing, Security, Hacking, Linux},
  pdfcenterwindow=true,
  colorlinks=true,
  breaklinks=true,
  linkcolor=blue,
  menucolor=blue,
  urlcolor=blue
}

% syntax highlighting
% all options should be set here document wide
\lstset{
backgroundcolor=\color[HTML]{EEEEEE},
frame=single,
basicstyle=\footnotesize\ttfamily,
float,
deletekeywords={return},
otherkeywords={mkdir, curl, sudo, sha1sum, grep, cut, sort, wget, makepkg,
pacman, blackman, chmod},
keywordstyle=\color{orange},
commentstyle=\color{blue},
stringstyle=\color{red},
language=bash,
showspaces=false,
showtabs=false,
tabsize=2
}

%%% HEADER / FOOTER %%%
\setlength{\headheight}{33pt}
\setlength{\headsep}{33pt}
\lhead{{\includegraphics[width=1cm,height=1cm]{images/logo.png}}}
\rhead{The BlackArch Linux Guide}

%%% CUSTOM MACROS %%%
% for html links
\ifpdf\else
\def\href#1#2{\htmladdnormallink{#2}{#1}}
\fi

%------------------%
%  TITLE PAGE      %
%------------------%
\begin{document}
\pagestyle{empty}
\begin{center}
\begin{figure}[htbp]
\centering
\vspace{0.5cm}
\includegraphics[width=8cm]{images/logo.png}
\label{fig:logo}
\end{figure}
\vspace{0.5cm}
\Huge{\textbf{The BlackArch Linux Guide}}\\
\vspace{1cm}
\Large{\color{blue}https://www.blackarch.org/}\\
\vspace{0.5cm}
\end{center}
\newpage
\tableofcontents
\newpage
\pagestyle{fancy}

%------------------%
%  Chapter 1       %
%------------------%

\chapter{Introducere}

\section{Overview}
Manualul pentru BlackArch Linux este impartit in cateva parti:
\begin{itemize}

\item Introducere-Ofera o previzualizare, o introducere, niste informatii in plus
\item Manualul utilizatorului-Totul ce un utilizator trebuie sa stie pentru a folosi BlackArch eficient
\item Manualul Dezvoltatorului-Cum sa contribui in proiect
\item Manualul pentru ustensile-Explicatii clare si precise despre folosirea lor (WIP)
\end{itemize}

\section{Ce este BlackArch Linux?}
BlackArch este o distributie Linux completa pentru securitate cibernetica

Este derivata de la  \href{https://www.archlinux.org/}{ArchLinux} iar utilizatorii pot instala componnentele BlackArch separat

Intregul sistem este distribuit ca Arch Linux \href{https://wiki.archlinux.org/index.php/Unofficial\_User\_Repositories}
{unofficial user repository} deci, poti instala BlackArch intr-o instalatie ArchLinux obisnuita.Pachetele pot fi instalate separat sau pe categorii

Este continuu actualizat continand peste
\href{https://www.blackarch.org/tools.html}{2600} unelte.
Toate uneltele sunt testate riguros  si adaugate intr-o baza de date pentru a fi folosite la capacitate maxima
% should quickly describe the testing methods/code review procedures etc

\section{Istoria BlackArch Linux}
Incurand...

\section{Platforme Suportate}
Incurand...

\section{Ajuta-ne}
Poti contacta echipa BlackArch print urmatoarele metode:

Website: \url{https://www.blackarch.org/}

Mail: \href{mailto:team@blackarch.org}{team@blackarch.org}

IRC: \url{irc://irc.freenode.net/blackarch}

Twitter: \url{https://twitter.com/blackarchlinux}

Github: \url{https://github.com/Blackarch/}

Discord: \url{https://discord.com/invite/xMHt8dW}

%------------------%
%  Chapter 2       %
%------------------%


\chapter{Manualul Utilizatorului}

\section{Instalare}

Urmatoarea sectiune iti va arata cum sa instalezi  BlackArch repository si cum sa instalezi pachete. BlackArch suporta ambele variante, instalarea din repository folosind
pachete binare cat si compiland si instaland din sursa.

BlackArch este compatibil cu normala instalare Arch.Este exact ca un neoficial 'pachet'. Daca vrei un ISO, mai jos
\href{https://www.blackarch.org/downloads.html#iso}{Live ISO} section.

\subsection{Instalarea sa pe ArchLinux}
Ruleaza \href{https://blackarch.org/strap.sh}{strap.sh} ca root urmand urmatoarele instructiuni. Urmatorul exemplu
\begin{lstlisting}
   curl -O https://blackarch.org/strap.sh
   sha1sum strap.sh # should match: d062038042c5f141755ea39dbd615e6ff9e23121
   sudo chmod +x strap.sh
   sudo ./strap.sh
\end{lstlisting}

Acum istaleaza o copie noua a pachetului master listeaza si sincronizeaza pachetele:
\begin{lstlisting}
  sudo pacman -Syyu
\end{lstlisting}



\subsection{Instalarea pachetelor}
Acum poti instala uneltele din blackarch repository.
\begin{enumerate}
\item Pentru a lista uneltele valabile, ruleaza
\begin{lstlisting}
  pacman -Sgg | grep blackarch | cut -d' ' -f2 | sort -u
\end{lstlisting}

\item Petru a instala toate uneltele, ruleaza
\begin{lstlisting}
  pacman -S blackarch
\end{lstlisting}

\item Pentru a instala o anumita categorie, ruleaza
\begin{lstlisting}
  pacman -S blackarch-<category>
\end{lstlisting}

\item Pentru a instala categoriile BlackArch, ruleaza
\begin{lstlisting}
  pacman -Sg | grep blackarch
\end{lstlisting}

\end{enumerate}


\subsection{Instalarea pachetelor din sursa}
Ca parte din o, poti construi pachete  BlackArch din sursa. Poti gasi PKGBUILDs in
\href{https://github.com/BlackArch/blackarch/tree/master/packages}{github}. To
creeaza intregul repo, poti folosi
\href{https://github.com/BlackArch/blackman}{Blackman} tool.
\begin{itemize}
\item Prima data, trebuie sa instalezi Blackman. Daca repo-ul pentru pachetul BlackArch
este configurat, poti instala Blackman:
\begin{lstlisting}
  pacman -S blackman
\end{lstlisting}

\item Poti instala si configura Blackman din sursa:
\begin{lstlisting}
  mkdir blackman
  cd blackman
  wget https://raw.github.com/BlackArch/blackarch/master/packages/blackman/PKGBUILD
  # Make sure the PKGBUILD has not been maliciously tampered with.
  makepkg -s
\end{lstlisting}

\item Sau poti instala BlackArch din AUR:
\begin{lstlisting}
  <whatever AUR helper you use> -S blackman
\end{lstlisting}

\end{itemize}
\subsection{Folosirea simpla a lui Blackman} BlackMan e foarte usor de folosit, desi e diferent de ce te-ai fi asteptat, e diferit de pacman.
\begin{itemize}
\item Downloadeaza, compileaza si instaleaza pachete:
\begin{lstlisting}
  sudo blackman -i package
\end{lstlisting}

\item Descarca, compileaza si instaleaza intreaga categorie:
\begin{lstlisting}
  sudo blackman -g group
\end{lstlisting}

\item Descarca, compileaza si instaleaza toate uneltele BlackArch:
\begin{lstlisting}
  sudo blackman -a
\end{lstlisting}

\item Pentru a lista categoriile BlackArch:
\begin{lstlisting}
  blackman -l
\end{lstlisting}

\item Pentru a lista categoriile uneltelor:
\begin{lstlisting}
  blackman -p category
\end{lstlisting}

\end{itemize}

\subsection{Instalarea din  live-, netinstall- ISO sau ArchLinux}
Poti instala BlackArch Linux din unul dintre live- neinstall-.ISO.\\See
\url{https://www.blackarch.org/download.html#iso}. Urmatorii pasi sunt necesari dupa boot up-ul ISO-ului
\begin{itemize}
\item Instalarea instalarului-blackarch:
\begin{lstlisting}
  sudo pacman -S blackarch-installer
\end{lstlisting}

\item Run
\begin{lstlisting}
  sudo blackarch-install
\end{lstlisting}

\end{itemize}

%------------------%
%  Chapter 3       %
%------------------%

\chapter{Manualul creatorului}

\section{Sistemul de construire al Arch si pachete separate}


PKGBUILD sunt fisiere. Fiecare ii  comunica  makepkg(1) cum sa creeze un
pachet. Fisierele PKGBUILD  sunt scrise in Bash.

Pentru mai multe informatii, citeste (sau rasfoieste) urmatoarele:
\begin{itemize}
\item \href{https://wiki.archlinux.org/index.php/Creating_Packages}{Arch Wiki: Crearea pachetelor}
\item \href{https://wiki.archlinux.org/index.php/Makepkg}{Arch Wiki: makepkg}
\item \href{https://wiki.archlinux.org/index.php/PKGBUILD}{Arch Wiki: PKGBUILD}
\item \href{https://wiki.archlinux.org/index.php/Arch_Packaging_Standards}{Arch Wiki:  inpacetarea standard a Arch }
\end{itemize}

\section{PKGBUILD pentru standardele BlackArch}
Pentru o mai buna intelegere,  PKGBUILDs sunt similare cu cele AUR,
cu mici diferente enumerate mai jos. Fiecare pachet trebuie sa apartina BlackArch catusi de putin.

\subsection{Grupuri}
Pentru a instala o specifica ramura de pachete rapid si usor,pachetele au fost
separate in grupuri. Acestea se pot instala prin simplul 'pacman -S <group name>' pentru a instala multiple pachete

\subsubsection{blackarch}
Grupul BlackArch este baza tuturor grupurilor is este locul in care pachetele thebuie sa vie stocate. Acest lucru da voie utilizatorilor sa instaleze pachete usor
Ce ar trebuie sa contina: TOTUL.

\subsubsection{blackarch-anti-forensic}
Pachetele ce sunt necesare pentru forensic,
incluzand encryption, steganography, si orice  ar trebuie sa modifice fisiere.
Aceste unelte ar trebui sa functioneze in cazul in care doriti sa ascundeti informatii.

Exemple: luks, TrueCrypt, Timestomp, dd, ropeadope, secure-delete

\subsubsection{blackarch-automation}
Pachete ce ar trebui folosite pentru workflow automation

Exemple: blueranger, tiger, wiffy

\subsubsection{blackarch-backdoor}
Pachete ce  exploateaza sau folosesc backdoors sunt  deja  in sisteme vulnerabile.

Exemple: backdoor-factory, rrs, weevely

\subsubsection{blackarch-binary}
Pachete ce opereaza pe fisiere binare.

Exemple: binwally, packerid

\subsubsection{blackarch-bluetooth}
Pachete ce ataca orice are legatura cu Bluetooth standard (802.15.1).

Exemple: ubertooth, tbear, redfang

\subsubsection{blackarch-code-audit}
Pachete ce  scaneaza o sursa existenta de cod in cautare de vulnerabilitati.

Exemple: flawfinder, pscan

\subsubsection{blackarch-cracker}
Pachete folosite pentru cracking cryptographic functions, ie hashes.

Exemple: hashcat, john, crunch

\subsubsection{blackarch-crypto}
Pachete ce folosesc  cryptography,cu exeptia decriptarii.

Exemple: ciphertest, xortool, sbd

\subsubsection{blackarch-database}
Packages that involve database exploitations on any level.

Examples: metacoretex, blindsql

\subsubsection{blackarch-debugger}
Pachete ce dau voie utilizatorului sa vizioneza ce ''face' un program intr-un anumit timp.

Exemple: radare2, shellnoob

\subsubsection{blackarch-decompiler}
Pachete ce incearca sa face  reverse la un program pentru ai vedea codul sursa.

Examples: flasm, jd-gui

\subsubsection{blackarch-defensive}
Pachete ce sunt folosite pentru a apara utilizatorul de virusi si atacuri.

Exemple: arpon, chkrootkit, sniffjoke

\subsubsection{blackarch-disassembler}
Este similar cu blackarch-decompiler, si vor fi probabil multe programe ce vor ajunge in mai multe catregorii, desi aceste pachete produc assembly output

Exemple: inguma, radare2

\subsubsection{blackarch-dos}
Pachete folosite pentru D0S  (Denial of Service).

Exemple: 42zip, nkiller2

\subsubsection{blackarch-drone}
Pachete ce sunt folosite pentru a folosi drone

Exemple: meshdeck, skyjack

\subsubsection{blackarch-exploitation}
Pachete ce se folosesc de vulnerabilitatiile altor programe sau servicii

Exemple: armitage, metasploit, zarp

\subsubsection{blackarch-fingerprint}
Pachete ce exploateaza  echipament fingerprint biometric.

Exemple: dns-map, p0f, httprint

\subsubsection{blackarch-firmware}
Pachete ce expluateaza vulnerabilitati in firmware:

Exemple: None yet, amend asap.

\subsubsection{blackarch-forensic}
Pachhete ce sunt folosite pentru a descoperi data pe disk-uri si memorie interna.

Exemple: aesfix, nfex, wyd

\subsubsection{blackarch-fuzzer}
Pachete ce sunt folosite pentru fuzz, ie
'aruncand' valori "la nimereala" in campuri de input pentru a vedea ce se intampla.

Exemple: msf, mdk3, wfuzz

\subsubsection{blackarch-hardware}
Pachete ce exploateaza sau folosesc orice are legatura cu hardware.

Exemple: arduino, smali

\subsubsection{blackarch-honeypot}
Pachete ce au legatura cu "honeypots", ie, programe ce au tendinta sa para vulnerabile pentru a fii exploatate.

Exemple: artillery, bluepot, wifi-honey

\subsubsection{blackarch-keylogger}
Pachetele ce tin cont de functionalitatea sistemului tau.

Exemple: None yet, amend asap.

\subsubsection{blackarch-malware}
Pachete ce tin cont de malware.

Exemple: malwaredetect, peepdf, yara

\subsubsection{blackarch-misc}
Pachte ce nu se incadreaza in vreo categorie.

Exemple: oh-my-zsh-git, winexe, stompy

\subsubsection{blackarch-mobile}
Pachete ce manipuleaza platforme mobile.

Exemple: android-sdk-platform-tools, android-udev-rules

\subsubsection{blackarch-networking}
Pachetel ce au legatura cu ip-ul.

Exemple: Cam toate 

\subsubsection{blackarch-nfc}
Pachete  ce se folosesc ce nfc(near-field communications).

Exemple: nfcutils

\subsubsection{blackarch-packer}
Pachete ce se folesesc de alte pachete.

\textit{packers sunt programe ce contin malware.}

Exemple: packerid

\subsubsection{blackarch-proxy}
Pachete ce sunt considerate ca proxy, ie redirectionand date /trafic
print internet.

Exemple: burpsuite, ratproxy, sslnuke

\subsubsection{blackarch-recon}
Pachete ce cauta vulnerabilitati. Mai mult o adunare decat un grup.

Exemple: canri, dnsrecon, netmask

\subsubsection{blackarch-reversing}
O adunare pentru dezasamblare

Exemple: capstone, radare2, zerowine

\subsubsection{blackarch-scanner}
Pachete ce cauta vulnerabilitati in sisteme specifice .

Exemple: scanssh, tiger, zmap

\subsubsection{blackarch-sniffer}
Pachete ce se concentreaza pe analizat traficul de date.

Exemple: hexinject, pytactle, xspy

\subsubsection{blackarch-social}
Pachete pentru social networking sites.

Exemple: jigsaw, websploit

\subsubsection{blackarch-spoof}
Pachete ce incearca sa faca spoof atacatorul, unde atacatorul devine o victima

Exemple: arpoison, lans, netcommander

\subsubsection{blackarch-threat-model}
Pachete  ce se folosesc pentru reporting/recording o amenintare.

Exemple: magictree

\subsubsection{blackarch-tunnel}
Pachete ce sunt folosite pentru a directiona network traffic la un anumit network.

Exemple : ctunnel, iodine, ptunnel

\subsubsection{blackarch-unpacker}
Pachete ce sunt  folosite pentru a extracta pre-facut malware din un
executabil.

Exemple: js-beautify

\subsubsection{blackarch-voip}
Programe ce opereaza pe voip programs si protocoale.

Exemple: iaxflood, rtp-flood, teardown

\subsubsection{blackarch-webapp}
Programe ce opereaza pe internet-facing applications.

Exemple: metoscan, whatweb, zaproxy

\subsubsection{blackarch-windows}
Aceasta categorie este alcatuita de programe ce ruleaza nativ pe windows,  iar noi folosim wine pentru a le folosi pe BlackArch.

Exemple: 3proxy-win32, pwdump, winexe

\subsubsection{blackarch-wireless}
Programe ce opereaza pe wireless networks la orice nivel.

Exemple: airpwn, mdk3, wiffy

\section{Repository structure}
Poti gasi intregul BlackArch repo aici:
\href{https://github.com/BlackArch/blackarch}{https://github.com/BlackArch/blackarch}.
Avem is cateva repos secundare aici:
\href{https://github.com/BlackArch}{https://github.com/BlackArch}.

In repo-ul principal de pe git, exista trei categorii imprtante:

\begin{itemize}
\item docs - Documentatie.
\item packages - fisiere PKGBUILD.
\item scripts - Scripturi folositoare.
\end{itemize}

\subsection{Scripts}
O referinta pentru scripturi se gaseste aici  \verb|scripts/| directory:

\begin{itemize}
\item baaur-Incurand, acesta va uploada fisiere pe  AUR.
\item babuild-Va construi un pachet.
\item bachroot-Va alege un chroot pentru folosire.
\item baclean-Va curata pachete vechi ca.pkg.tar.xz din  repo.
\item baconflict-Incurand va inlocui \verb|scripts/conflicts|.
\item bad-files-Va gasi fisiere corupte.
\item balock-Va detine sau dona pachete.
\item banotify-Va notifica IRC despre package pushes.
\item barelease-Va da pachete pe repo.
\item baright-Ofera  BlackArch copyright info.
\item basign-Pachete de logare.
\item basign-key-Va asocia o cheie.
\item blackman-Are acelasi principiu ca pacman dar pentru git(a nu se confunda cu nrz's Blackman).
\item check-groups-Verifica grupurile.
\item checkpkgs-Verifica pachete pentru erori.
\item conflicts-Verifica pentru conflicturi de pachete.
\item dbmod-Modifica baza de date pentru pachete.
\item depth-list-Creaza o lista pentru dependency depth.
\item deptree-Creeaza un dependency tree,listeaza doar pachete blackarch.
\item get-blackarch-deps-Ia o lista pentru blackarch dependencies pentru un pachet.
\item get-official-Ia pachete oficiale.
\item list-loose-packages-Listeaza pachete ce nu sunt in grupurisi nu au dependente de alte pachete.
\item list-needed-Listeaza dependente lipsa.
\item list-removed-Listeaza pachete ce sunt in  repo dar nu in git.
\item list-tools-Listeaza  unelte.
\item outdated-Verifica pachete out-dated in pachetul repo ca cel de pe git.
\item pkgmod-Modifica un pachet creat.
\item pkgrel-Include pkgrel intr-un pachet.
\item prep-Curata un PKGBUILD al unui fisier si gaseste erori.
\item sitesync-Sync o copie locala a unui pachet repo si o copie remote.
\item size-hunt-Cauta pachete mari.
\item source-backup-Face backup la pachete.
\end{itemize}

\section{Contribuie la proiect}
In aceasta sectiune va fi prezentat cum sa participi la proiectul BlackArch Linux. Aceptam cereri pull de toate dimensiunile, de la unele minore la pachete intregi.\\Pentru ajutor sau intrebari ne puteti contacta.

\\\\
Toate lumea poate participa. Orice contributie e aceptata.

\subsection{Tutoriale recomandate}
Va rugam sa cititi urmatoarele inainte de a ne ajuta:
\begin{itemize}
\item
\href{https://wiki.archlinux.org/index.php/Arch\_Packaging\_Standards)}{Arch
Packaging Standards}
\item \href{https://wiki.archlinux.org/index.php/Creating\_Packages}{Creating
Packages}
\item \href{https://wiki.archlinux.org/index.php/PKGBUILD}{PKGBUILD}
\item \href{https://wiki.archlinux.org/index.php/Makepkg}{Makepkg}
\end{itemize}

\subsection{Pasii pentru a contribui}
 Pentru a ne trimite schimbariile pentru proiectul BlackArchLinux, urmati urmatorii pasi:

\begin{enumerate}
\item Fork the repository from
\url{https://github.com/BlackArch/blackarch}
\item Modificati fisierele de care aveti nevoie, (e.g. PKGBUILD, .patch files, etc).
\item Pregatiti schimbariile.
\item Da-ti push la schimbari.
\item Spuneti-ne sa va aceptam schimbarile, preferabil printr-un pull.
\end{enumerate}

\subsection{Exemple}
Urmatorul exemplu demonstreaza cum se va trimite un nou pachet catre BlackArch. Folosim \href{https://wiki.archlinux.org/index.php/yaourt}{yaourt}
(poti folosi pacaur deasemenea) pentru a lua un pachet pre-existent fisierul PKGBUILD pentru \textbf{nfsshell} din \href{https://aur.archlinux.org/}{AUR} si modifical cum vrei

\subsubsection{Foloseste PKGBUILD}
Ia \textit{PKGBUILD} fisierul folosind yaourt sau pacaur:
\begin{lstlisting}
  user@blackarchlinux $ yaourt -G nfsshell
  ==> Download nfsshell sources
  x LICENSE
  x PKGBUILD
  x gcc.patch
  user@blackarchlinux $ cd nfsshell/
\end{lstlisting}

\subsubsection{Curata PKGBUILD}
Curata \textit{PKGBUILD} fisierul si salveaza:
\begin{lstlisting}
  user@blackarchlinux nfsshell $ ./blackarch/scripts/prep PKGBUILD
  cleaning 'PKGBUILD'...
  expanding tabs...
  removing vim modeline...
  removing id comment...
  removing contributor and maintainer comments...
  squeezing extra blank lines...
  removing '|| return'...
  removing leading blank line...
  removing $pkgname...
  removing trailing whitespace...
\end{lstlisting}

\subsubsection{Adjust PKGBUILD}
Ajunstaza \textit{PKGBUILD} fisierul:
\begin{lstlisting}
  user@blackarchlinux nfsshell $ vi PKGBUILD
\end{lstlisting}

\subsubsection{Build the package}
Construieste pachetul:
\begin{lstlisting}user@blackarchlinux nfsshell $ makepkg -sf
==> Making package: nfsshell 19980519-1 (Mon Dec  2 17:23:51 CET 2013)
==> Checking runtime dependencies...
==> Checking buildtime dependencies...
==> Retrieving sources...
-> Downloading nfsshell.tar.gz...
% Total    % Received % Xferd  Average Speed   Time    Time     Time
CurrentDload  Upload   Total   Spent    Left  Speed100 29213  100 29213    0
0  48150      0 --:--:-- --:--:-- --:--:-- 48206
-> Found gcc.patch
-> Found LICENSE
...
<lots of build process and compiler output here>
...
==> Leaving fakeroot environment.
==> Finished making: nfsshell 19980519-1 (Mon Dec  2 17:23:53 CET 2013)
\end{lstlisting}

\subsubsection{Instaleaza si testeaza pachetul}
Instaleaza si testeaza pachetul:
\begin{lstlisting}
  user@blackarchlinux nfsshell $ pacman -U nfsshell-19980519-1-x86_64.pkg.tar.xz
  user@blackarchlinux nfsshell $ nfsshell # test it
\end{lstlisting}

\subsubsection{Adauga, modifica si trimite pachetul}
Adauga, confirma si trimite pachetul
\begin{lstlisting}user@blackarchlinux nfsshell $ cd /blackarchlinux/packages
user@blackarchlinux ~/blackarchlinux/packages $ mv ~/nfsshell .
user@blackarchlinux ~/blackarchlinux/packages $ git commit -am nfsshell && git push
\end{lstlisting}

\subsubsection{Creeaza o cerere pull }
Creeaza o cerere pe \href{https://github.com/}{github.com}
\begin{lstlisting}
  firefox https://github.com/<contributor>/blackarchlinux
\end{lstlisting}

\subsubsection{Adauga un remote pentru upstream}
Un lucru istet de facut este sa dai fork si sa lucrezi remote .
\begin{lstlisting}
  user@blackarchlinux ~/blackarchlinux $ git remote -v
  origin <the url of your fork> (fetch)
  origin <the url of your fork> (push)
  user@blackarchlinux ~/blackarchlinux $ git remote add upstream https://github.com/blackarch/blackarch
  user@blackarchlinux ~/blackarchlinux $ git remote -v
  origin <the url of your fork> (fetch)
  origin <the url of your fork> (push)
  upstream https://github.com/blackarch/blackarch (fetch)
  upstream https://github.com/blackarch/blackarch (push)
\end{lstlisting}

In mod implicit, git ar trebui sa trimita direct la origine, dar ca sa te asiguri ca git config este configurat corect.Acest lucru nu ar trebui sa fie o problema decat daca trebuie sa trimiti drepturi si nu poti trimite upstream fara ele

Daca poti trimite, ai putea aveam mai mult success cu
\textit{git@github.com:blackarch/blackarch.git} dar e alegerea ta.

\subsection{Requests}
\begin{enumerate}
\item Nu adauga \textbf{Maintainer} sau \textbf{Contributor} comentari pentru fisierele
\textit{PKGBUILD} . Adauga utilizator pentru sectiune AUTHORS  BlackArch guide.
\item Pentru o mai buna intelegere, va rugam sa respectati  stilul tuturor
\textit{PKGBUILD} fisierelor din repo.
\end{enumerate}

\subsection{Recomandari}
\href{http://wiki.archlinux.org/index.php/Namcap}{namcap} poate verifica pentru erori.

%------------------%
%  Chapter 4       %
%------------------%

\chapter{Tools Guide}
In curand...

\section{Coming Soon}
In curand...

%%% APPENDIX %%%
\appendix
%%%%%%%%%%%%%%%%%%%%%%%%%%%%%%%%%%%%%%%%%%%%%%%%%%%%%%%%%%%%%%%%%%%%%%%%%%%%%%%%
%                                                                              %
% BlackArch Linux appendix template                                            %
%                                                                              %
%%%%%%%%%%%%%%%%%%%%%%%%%%%%%%%%%%%%%%%%%%%%%%%%%%%%%%%%%%%%%%%%%%%%%%%%%%%%%%%%

\appendix

\chapter{Appendix}

\section{FAQs}

\section{AUTHORS}
The following people have contributed directly to BlackArch:
\begin{itemize}
\item Tyler Bennnett (tylerb@trix2voip.com)
\item fnord0 (fnord0@riseup.net)
\item nrz (nrz@nullsecurity.net)
\item Ellis Kenyo (elken.tdos@gmail.com)
\item CaledoniaProject (the.warl0ck.1989@gmail.com)
\item sudokode (sudokode@gmail.com)
\item Valentin Churavy (v.churavy@gmail.com)
\item Boy Sandy Gladies Arriezona (reno.esper@gmail.com)
\item Mathias Nyman
\item Johannes Löthberg (demizide@gmail.com)
\item Thiago da Silva Teixeira (teixeira.zeus@gmail.com)
\end{itemize}
The following people contributed directly to ArchPwn, which has been merged
into BlackArch:
\begin{itemize}
\item Francesco Piccinno (stack.box@gmail.com)
\item jensp (jens@jenux.homelinux.org)
\item Valentin Churavy (v.churavy@gmail.com)
\end{itemize}
We have taken build code from the following people:
\begin{itemize}
\item  3ED (krzysztof1987@gmail.com)
\item  AUR Perl (aurperl@juster.info)
\item  Aaron Griffin (aaron@archlinux.org)
\item  Abakus (java5@arcor.de)
\item  Adam Wolk (netprobe@gmail.com)
\item  Aleix Pol (aleixpol@kde.org)
\item  Aleshus (aleshusi@gmail.com)
\item  Alessandro Pazzaglia (jackdroido@gmail.com)
\item  Alessandro Sagratini (ale\_sagra@hotmail.com)
\item  Alex Cartwright (alexc223@googlemail.com)
\item  Alexander De Sousa (archaur.xandy21@spamgourmet.com)
\item  Alexander Rødseth (rodseth@gmail.com)
\item  Allan McRae (allan@archlinux.org)
\item  AmaN (gabroo.punjab.da@gmail.com)
\item  Andre Klitzing (aklitzing@online.de)
\item  Andrea Scarpino (andrea@archlinux.org)
\item  Andreas Schönfelder (passtschu@freenet.de)
\item  Andrej Gelenberg (andrej.gelenberg@udo.edu)
\item  Angel Velasquez (angvp@archlinux.org)
\item  Antoine Lubineau (antoine@lubignon.info)
\item  Anton Bazhenov (anton.bazhenov@gmail.com)
\item  Arkham (arkham@archlinux.us)
\item  Arthur Danskin (arthurdanskin@gmail.com)
\item  Balda (balda@balda.ch)
\item  Balló György (ballogyor+arch@gmail.com)
\item  Bartek Piotrowski (barthalion@gmail.com)
\item  Bartosz Feński (fenio@debian.org)
\item  Bartłomiej Piotrowski (nospam@bpiotrowski.pl)
\item  Bogdan Szczurek (thebodzio@gmail.com)
\item  Brad Fanella (bradfanella@archlinux.us)
\item  Brian Bidulock (bidulock@openss7.org)
\item  C Anthony Risinger (anthony@xtfx.me)
\item  CRT (crt.011@gmail.com)
\item  Can Celasun (dcelasun@gmail.com)
\item  Chaniyth (chaniyth@yahoo.com)
\item  Chris Brannon (cmbrannon79@gmail.com)
\item  Chris Giles (Chris.G.27@gmail.com) \& daschu117
\item  Christoph Siegenthaler (csi@gmx.ch)
\item  Christoph Zeiler (archNOSPAM@moonblade.org)
\item  Clément DEMOULINS (clement@archivel.fr)
\item  Corrado Primier (bardo@aur.archlinux.org)
\item  Daenyth (Daenyth+Arch@gmail.com)
\item  Dale Blount (dale@archlinux.org)
\item  Damir Perisa (damir.perisa@bluewin.ch)
\item  Dan Fuhry (dan@fuhry.us)
\item  Dan Serban (dserban01@yahoo.com)
\item  Daniel A. Campoverde Carrión
\item  Daniel Golle
\item  Daniel Griffiths (ghost1227@archlinux.us)
\item  Daniel J Griffiths (ghost1227@archlinux.us)
\item  Daniel Micay (danielmicay@gmail.com)
\item  Dave Reisner (dreisner@archlinux.org)
\item  Dawid Wrobel (cromo@klej.net)
\item  Devaev Maxim (mdevaev@gmail.com)
\item  Devin Cofer (ranguvar@archlinux.us)
\item  DigitalPathogen (aur@InfoSecResearchLabs.co.uk)
\item  DigitalPathogen (aur@digitalpathogen.co.uk)
\item  Dmitry A. Ilyashevich (dmitry.ilyashevich@gmail.com)
\item  Dominik Heidler (dheidler@gmail.com)
\item  DrZaius (lou@fakeoutdoorsman.com)
\item  Ebubekir KARUL (ebubekirkarul@yandex.com)
\item  Eduard "bekks" Warkentin (eduard.warkentin@gmail.com)
\item  Elmo Todurov (todurov@gmail.com)
\item  Emmanuel Gil Peyrot (linkmauve@linkmauve.fr)
\item  Eric Belanger (eric@archlinux.org)
\item  Ermak (ermak@email.it)
\item  Evangelos Foutras (evangelos@foutrelis.com)
\item  Fabian Melters (melters@gmail.com)
\item  Fabiano Furtado (fusca14@gmail.com)
\item  Federico Quagliata (ntp@quaqo.org)
\item  Firmicus (francois.archlinux@org)
\item  Florian Pritz (bluewind@jabber.ccc.de)
\item  Florian Pritz (flo@xinu.at)
\item  Francesco Piccinno (stack.box@gmail.com)
\item  François Charette (francois@archlinux.org)
\item  Gaetan Bisson (bisson@archlinux.org)
\item  Geoffroy Carrier (geoffroy.carrier@koon.fr)
\item  Georg Grabler (STiAT)
\item  George Hilliard (gh403@msstate.edu)
\item  Gerardo Exequiel Pozzi (vmlinuz386@yahoo.com.ar)
\item  Gilles CHAUVIN (gcnweb@gmail.com)
\item  Giovanni Scafora (giovanni@archlinux.org)
\item  Gordin (9ordin@gmail.com)
\item  Guillaume ALAUX (guillaume@archlinux.org)
\item  Guillermo Ramos (0xwille@gmail.com)
\item  Gustavo Alvarez (sl1pkn07@gmail.com)
\item  Hugo Doria (hugo@archlinux.org)
\item  Hyacinthe Cartiaux (hyacinthe.cartiaux@free.fr)
\item  James Fryman (jfryman@gmail.com)
\item  Jan "heftig" Steffens (jan.steffens@gmail.com)
\item  Jan de Groot (jgc@archlinux.org)
\item  Jaroslav Lichtblau (dragonlord@aur.archlinux.org)
\item  Jaroslaw Swierczynski (swiergot@aur.archlinux.org)
\item  Jason Chu (jason@archlinux.org)
\item  Jason R Begley (jayray@digitalgoat.com)
\item  Jason Rodriguez
\item  Jason St. John (jstjohn@purdue.edu)
\item  Jawmare (victor2008@gmail.com)
\item  Jeff Mickey (jeff@archlinux.org)
\item  Jens Pranaitis (jens@chaox.net)
\item  Jens Pranaitis (jens@jenux.homelinux.org)
\item  Jinx (jinxware@gmail.com)
\item  John D Jones III (jnbek1972@gmail.com)
\item  John Proctor (jproctor@prium.net)
\item  Jon Bergli Heier (snakebite@jvnv.net)
\item  Jonas Heinrich
\item  Jonathan Steel (jsteel@aur.archlinux.org)
\item  Joris Steyn (jorissteyn@gmail.com)
\item  Josh VanderLinden (arch@cloudlery.com)
\item  Jozef Riha (jose1711@gmail.com)
\item  Judd Vinet (jvinet@zeroflux.org)
\item  Juergen Hoetzel (jason@archlinux.org)
\item  Juergen Hoetzel (juergen@archlinux.org)
\item  Justin Davis (jrcd83@gmail.com)
\item  Kaiting Chen (kaitocracy@gmail.com)
\item  Kaos
\item  Kevin Piche (kevin@archlinux.org)
\item  Kory Woods (kory@virlo.net)
\item  Kyle Keen (keenerd@gmail.com)
\item  Larry Hajali (larryhaja@gmail.com)
\item  LeCrayonVert
\item  Le\_suisse (lesuisse.dev+aur@gmail.com)
\item  Lekensteyn (lekensteyn@gmail.com)
\item  Limao Luo (luolimao+AUR@gmail.com)
\item  Lucien Immink
\item  Lukas Fleischer (archlinux@cryptocrack.de)
\item  Manolis Tzanidakis
\item  Marcin "avalan" Falkiewicz (avalatron@gmail.com)
\item  Mariano Verdu (verdumariano@gmail.com)
\item  Marti Raudsepp (marti@juffo.org)
\item  MatToufoutu (mattoufootu@gmail.com)
\item  Matthew Sharpe (matt.sharpe@gmail.com)
\item  Mauro Andreolini (mauro.andreolini@unimore.it)
\item  Max Pray a.k.a. Synthead (synthead@gmail.com)
\item  Max Roder (maxroder@web.de)
\item  Maxwell Pray a.k.a. Synthead (synthead@gmail.com)
\item  Maxwell Pray a.k.a. Synthead (synthead1@gmail.com)
\item  Mech (tiago.bmp@gmail.com)
\item  Michael Düll (mail@akurei.me)
\item  Michael P (ptchinster@archlinux.us)
\item  Michal Krenek (mikos@sg1.cz)
\item  Michal Zalewski (lcamtuf@coredump.cx)
\item  Miguel Paolino (mpaolino@gmail.com)
\item  Miguel Revilla (yo@miguelrevilla.com)
\item  Mike Roberts (noodlesgc@gmail.com)
\item  Mike Sampson (mike@sambodata.com)
\item  Nassim Kacha (nassim.kacha@gmail.com)
\item  Nicolas Pouillard (nicolas.pouillard@gmail.com)
\item  Nicolas Pouillard https://nicolaspouillard.fr
\item  Niklas Schmuecker
\item  Oleander Reis (oleander@oleander.cc)
\item  Olivier Le Moal (mail@olivierlemoal.fr)
\item  Olivier Médoc "oliv" (o\_medoc@yahoo.fr)
\item  Pascal E. (archlinux@hardfalcon.net)
\item  Patrick Leslie Polzer (leslie.polzer@gmx.net)
\item  Paul Mattal (paul@archlinux.org)
\item  Paul Mattal (pjmattal@elys.com)
\item  Pengyu CHEN (cpy.prefers.you@gmail.com)
\item  Peter Wu (lekensteyn@gmail.com)
\item  Philipp 'TamCore' B. (philipp@tamcore.eu)
\item  Pierre Schmitz (pierre@archlinux.de)
\item  Pranay Kanwar (pranay.kanwar@gmail.com)
\item  Pranay Kanwar (warl0ck@metaeye.org)
\item  PyroPeter (googlemail@com.abi1789)
\item  PyroPeter (googlemail.com@abi1789)
\item  Ray Rashif (schiv@archlinux.org)
\item  Remi Gacogne
\item  Renan Fernandes (renan@kauamanga.com)
\item  Richard Murri (admin@richardmurri.com)
\item  Roberto Alsina (ralsina@kde.org)
\item  Robson Peixoto (robsonpeixoto@gmail.com)
\item  Roel Blaauwgeers (roel@ttys0.nl)
\item  Rorschach (r0rschach@lavabit.com)
\item  Ruben Schuller (shiml@orgizm.net)
\item  Rudy Matela (rudy@matela.com)
\item  Ryon Sherman (ryon.sherman@gmail.com)
\item  Sabart Otto \item  Seberm (seberm@gmail.com)
\item  SakalisC (chrissakalis@gmail.com)
\item  Sam Stuewe (halosghost@archlinux.info)
\item  SanskritFritz (SanskritFritz@gmail.com)
\item  Sarah Hay (sarahhay@mb.sympatico)
\item  Sebastian Benvenuti (sebastianbenvenuti@gmail.com)
\item  Sebastian Nowicki (sebnow@gmail.com)
\item  Sebastien Duquette (ekse.0x@gmail.com)
\item  Sebastien LEDUC (sebastien@sleduc.fr)
\item  Sebastien Leduc (sebastien@sleduc.fr)
\item  Sergej Pupykin (pupykin.s+arch@gmail.com)
\item  Sergio Rubio (rubiojr@biondofu.net)
\item  Sheng Yu (magicfish1990@gmail.com)
\item  Simon Busch (morphis@gravedo.de)
\item  Simon Legner (Simon.Legner@gmail.com)
\item  Sirat18 (aur@sirat18.de)
\item  SpepS (dreamspepser@yahoo.it)
\item  Spider.007 (archPackage@spider007.net)
\item  Stefan Seering
\item  Stephane Travostino (stephane.travostino@gmail.com)
\item  Stéphane Gaudreault (stephane@archlinux.org)
\item  Sven Kauber (celeon@gmail.com)
\item  Sven Schulz (omee@archlinux.de)
\item  Sébastien Duquette (ekse.0x@gmail.com)
\item  Sébastien Luttringer (seblu@archlinux.org)
\item  TDY (tdy@gmx.com)
\item  Teemu Rytilahti (tpr@iki.fi)
\item  Testuser\_01
\item  Thanx (thanxm@gmail.com)
\item  Thayer Williams (thayer@archlinux.org)
\item  Thomas S Hatch (thatch45@gmail.com)
\item  Thorsten Töpper
\item  Tilmann Becker (tilmann.becker@freenet.de)
\item  Timothy Redaelli (timothy.redaelli@gmail.com)
\item  Timothée Ravier (tim@siosm.fr)
\item  Tino Reichardt
\item  Tobias Kieslich (tobias@justdreams.de)
\item  Tobias Powalowski (tpowa@archlinux.org)
\item  Tom K (tomk@runbox.com)
\item  Tom Newsom (Jeepster@gmx.co.uk)
\item  Tomas Lindquist Olsen (tomas@famolsen.dk)
\item  Travis Willard (travisw@wmpub.ca)
\item  Valentin Churavy (v.churavy@gmail.com)
\item  ViNS (gladiator@fastwebnet.it)
\item  Vlatko Kosturjak (kost@linux.hr)
\item  Wes Brown (wesbrown18@gmail.com)
\item  William Rea (sillywilly@gmail.com)
\item  Xavier Devlamynck (magicrhesus@ouranos.be)
\item  Xiao\item Long Chen (chenxiaolong@cxl.epac.to)
\item  aeolist (aeolist@hotmail.com)
\item  ality@pbrane.org
\item  astaroth (astaroth\_@web.de)
\item  bender02@archlinux.us
\item  billycongo (billycongo@gmail.com)
\item  bslackr (brendan@vastactive.com)
\item  cbreaker (cbreaker@tlen.pl)
\item  chimeracoder (dev@chimeracoder.net)
\item  damir (damir@archlinux.org)
\item  danitool
\item  darkapex (me@jailuthra.in)
\item  daronin
\item  dkaylor (dpkaylor@gmail.com)
\item  dobo (dobo90\_at\_gmail@com)
\item  dorphell (dorphell@archlinux.org)
\item  evr (evanroman@at.gmail)
\item  fnord0 (fnord0@riseup.net)
\item  fxbru (frxbru@gmail)
\item  hcar
\item  icarus (icarus.roaming@gmail.com)
\item  iceman (icemanf@gmail.com)
\item  kastor (kastor@fobos.org)
\item  kfgz (kfgz@interia.pl)
\item  linuxSEAT (linuxSEAT@gmail.com)
\item  m4xm4n (max@maxfierke.com)
\item  mar77i (mysatyre@gmail.com)
\item  marc0s (marc0s@fsfe.org)
\item  mickael9 (mickael9@gmail.com)
\item  nblock (nblock@archlinux.us)
\item  nofxx (x@nofxx.com)
\item  onny (onny@project
\item  pootzko (pootzko@gmail.com)
\item  revel (revel@muub.net)
\item  rich\_o (rich\_o@lavabit.com)
\item  s1gma (s1gma@mindslicer.com)
\item  sandman (r.coded@gmail.com)
\item  sebikul (sebikul@gmail.com)
\item  sh0 (mee@sh0.org)
\item  shild (sxp@bk.ru)
\item  simo (simo@archlinux.org)
\item  snuo
\item  sudokode (sudokode@gmail.com)
\item  tobias (tobias@archlinux.org)
\item  trashstar (trash@ps3zone.org)
\item  unexist (unexist@subforge.org)
\item  untitled (rnd0x00@gmail.com)
\item  virtuemood (virtue@archlinux.us)
\item  wido (widomaker2k7@gmail.com)
\item  wodim (neikokz@gmail.com)
\item  yannsen (ynnsen@gmail.com)
\end{itemize}


\end{document}

%%% EOF %%%
